\chapter{Depth, vanishing, and nonvanishing}
\label{ch:depth-and-vanishing}

Let \(R\) be Noetherian and let \(M\) be finite. We use
\(\dim M=\dim\operatorname{Supp} M\) with \(\dim 0=-\infty\).
We use the depth and regular-sequence facts proved in
\cite[DP-01, DP-04--DP-06, CM-01--CM-04]{core}. In particular
\(\operatorname{grade}(I,M)\) is the length of a maximal
\(M\)-regular sequence in \(I\) when \(V(I)\cap\operatorname{Supp}M\)
is nonempty, with value \(+\infty\) in the disjoint case.

\begin{theorem}[Depth detected by local cohomology]\label{res:V01}
With \(\inf\varnothing=+\infty\),
\[
 \operatorname{grade}(I,M)=\inf\{q\geq0:H_I^q(M)\ne0\}.
\]
\end{theorem}
\begin{proof}
If the supports are disjoint, Result~\ref{res:L08} gives zero in every
degree. Otherwise the grade is finite by the local depth formula of
\cite[DP-06]{core}. At grade zero there is an associated prime of \(M\)
containing \(I\), which supplies a nonzero element killed by \(I\).
For positive grade \(g\), prime avoidance supplies a regular
\(x\in I\) and \(M/xM\) has grade \(g-1\). It has nonempty support
in \(V(I)\) by local Nakayama at any prime of the original intersection.
Apply induction and the long exact sequence for
\(0\to M\xrightarrow{x}M\to M/xM\to0\).
For \(q<g\), multiplication by \(x\) on \(H_I^q(M)\) is injective
because \(H_I^{q-1}(M/xM)=0\). Since this module is \(I\)-power
torsion by Result~\ref{res:L02}, it is zero. At degree \(g\) the
nonzero group \(H_I^{g-1}(M/xM)\) injects into \(H_I^g(M)\).
\end{proof}

\begin{theorem}[Grothendieck vanishing]\label{res:V02}
If \(M\) has finite-dimensional support, \(H_I^q(M)=0\) for
\(q>\dim M\).
\end{theorem}
\begin{proof}
A finite module over a Noetherian ring has a finite filtration
with factors \(R/\mathfrak p\). To construct one, choose in each
nonzero quotient a cyclic submodule whose annihilator is a maximal
annihilator of a nonzero element. That annihilator is prime, since
if \(ab\) annihilates the generator and \(b\) does not, the
annihilator of its multiple by \(b\) contains \(a\) and cannot
strictly enlarge the original one. Noetherianity terminates the
resulting chain. All factors have support inside that of \(M\).
The long exact sequence and Result~\ref{res:L08} reduce the claim
to a domain \(B\) of dimension at most \(\dim M\).
Zero and unit ideals follow from Result~\ref{res:L06}. Otherwise take
\(0\ne x\in IB\). Prepending the zero prime to any chain in \(V(x)\)
shows \(\dim B/xB\leq\dim B-1\). Induct on dimension, applying the
claim for lower-dimensional modules to \(B/xB\).
The sequence for multiplication by \(x\) makes
\(x:H^q_{IB}(B)\to H^q_{IB}(B)\) injective for \(q>\dim B\).
Power torsion then gives zero. At dimension zero, a domain is a field
and there are only the two extreme ideals.
\end{proof}

\begin{lemma}[A regular-local top class]\label{res:V03}
If \(S\) is regular local of dimension \(d\) and
\(x_1,\ldots,x_d\) minimally generate its maximal ideal \(\mathfrak n\),
the class of \(1/(x_1\cdots x_d)\) in \(H^d_{\mathfrak n}(S)\)
is nonzero. Multiplication by each \(0\ne c\in S\) on this
cohomology module is surjective.
\end{lemma}
\begin{proof}
The minimal generators are regular and permutable and their positive
powers are likewise regular in any order
(\cite[ALG-05-C03-S03]{p103} and \cite[RK-09, CM-08]{core}).
For each \(n\), take the colon of the ideal
\((x_1^{n+1},\ldots,x_d^{n+1})\) by the product \(x_1\cdots x_d\),
one factor at a time. At a factor \(x_j\), reduce by the powers of
the other parameters. In the quotient \(x_j\) is a nonzerodivisor,
so taking the colon reduces its exponent by one. Thus the final
colon is \((x_1^n,\ldots,x_d^n)\).
The maps of the top colimit in Result~\ref{res:L07} are injective,
so the class of \(1\) at stage one survives. A regular local ring is
a domain. If \(c\) is a nonunit, \(S/cS\) has smaller dimension by
the prime-chain argument of Result~\ref{res:V02}. Its \(d\)-th local
cohomology vanishes, so the long exact sequence proves surjectivity.
For a unit multiplication is an isomorphism. When \(d=0\), \(S\)
is a field and the class is \(1\) in \(H^0_{(0)}(S)=S\).
\end{proof}

\begin{proposition}[Completion preserves local depth and dimension]\label{res:V04}
Let \(R\) be local and \(M\ne0\) finite. Then \(\widehat M\) has
the same dimension and depth as \(M\), and the \(\mathfrak m\)-local
cohomology groups of \(M\) identify with the
\(\mathfrak m\widehat R\)-local groups of \(\widehat M\).
\end{proposition}
\begin{proof}
The last assertion is Result~\ref{res:L09} and the tensor formula for
completion. Faithful flatness ensures \(\widehat M\ne0\).
The Hilbert--Samuel functions of \(M\) and \(\widehat M\) agree
because \(M/\mathfrak m^{n+1}M\) and
\(\widehat M/(\mathfrak m\widehat R)^{n+1}\widehat M\) agree over
the same Artin quotient, so their composition lengths agree.
Their degrees are the support dimensions by
\cite[ALG-05-C03-S02, ALG-05-C04-R03 and S01]{p103}.
Apply Result~\ref{res:V01} at the maximal ideals for depth.
\end{proof}

\begin{lemma}[Top cohomology in positive rank]\label{res:V05}
Let \(B\) be a complete Noetherian local domain of dimension \(d\).
If a finite \(B\)-module \(N\) has positive rank, then
\(H^d_{\mathfrak m_B}(N)\ne0\).
\end{lemma}
\begin{proof}
By Result~\ref{res:C05}, \(B\) is finite over a regular local
subring \((S,\mathfrak n)\) of dimension \(d\), with
\(\sqrt{\mathfrak nB}=\mathfrak m_B\).
The fraction field of \(B\) is a finite extension of that of \(S\).
localization of \(B\) at the nonzero elements of \(S\) is a finite domain
over a field, hence a field. Thus \(N\) has positive rank over \(S\).
A nonzero linear functional over the fraction field of \(S\) can be
scaled by a denominator common to a finite set of generators to give
a nonzero \(S\)-linear map \(N\to S\). Choose \(u\in N\) whose image
\(c\) is nonzero. The composite \(S\to N\to S\), where the first
arrow sends \(1\) to \(u\), is multiplication by \(c\).
The induced composite on \(H^d_{\mathfrak n}(S)\) is nonzero by
Result~\ref{res:V03}. It factors through \(H^d_{\mathfrak n}(N)\),
which identifies with the claimed module by Results~\ref{res:L06}
and~\ref{res:L08}.
\end{proof}

\begin{theorem}[Grothendieck nonvanishing]\label{res:V06}
For a nonzero finite module \(M\) over a Noetherian local ring,
\(H_{\mathfrak m}^{\dim M}(M)\ne0\).
\end{theorem}
\begin{proof}
By Result~\ref{res:V04}, we may work over the completion with the
same dimension \(d\). Choose \(\mathfrak p\) in the support with
\(\dim R/\mathfrak p=d\), by taking the bottom of a chain of maximal
length. Put \(N=M/\mathfrak pM\). Its localization at the generic
point of the complete domain \(B=R/\mathfrak p\) is nonzero by
Nakayama over \(R_{\mathfrak p}\). Hence \(N\) has positive rank over
\(B\). The long exact sequence and Result~\ref{res:V02} give a
surjection \(H_{\mathfrak m}^d(M)\to H_{\mathfrak m}^d(N)\), since
\(\mathfrak p M\) is supported in \(\operatorname{Supp}M\).
The target is nonzero by Results~\ref{res:L08} and~\ref{res:V05}.
\end{proof}

\begin{corollary}[Cohen--Macaulay criterion]\label{res:V07}
A nonzero finite module over a Noetherian local ring is
Cohen--Macaulay if and only if \(H^q_{\mathfrak m}(M)\) is
concentrated in degree \(\dim M\). Its group in that degree is nonzero.
\end{corollary}
\begin{proof}
By Results~\ref{res:V01}, \ref{res:V02} and~\ref{res:V06},
the bottom and top nonzero degrees are respectively depth and
dimension. Equality is the definition of Cohen--Macaulayness.
The assertion for the ring follows by taking \(M=R\).
\end{proof}

\begin{corollary}[Extension across the support]\label{res:V08}
For \(U=\operatorname{Spec}R\setminus V(I)\), the restriction
\(M\to\Gamma(U,\widetilde M)\) is an isomorphism if and only if
\(\operatorname{grade}(I,M)\geq2\).
\end{corollary}
\begin{proof}
Its kernel and cokernel are \(H_I^0(M)\) and \(H_I^1(M)\) by
Result~\ref{res:L10}. Both vanish precisely under the asserted
condition by Result~\ref{res:V01}.
\end{proof}
